\documentclass{stex}

\libusepackage{naproche}
\libinput{preamble}

\begin{document}
\begin{smodule}{natural-numbers.ftl}

\importmodule[libraries/foundations]{classes.ftl}

\symdef{Naturals}{\mathbb N}
\symdef{Plus}[name=sum]{\,+\,}
\symdef{Zero}[name=zero]{0}
\symdef{One}[name=one]{1}
\symdef{Two}[name=two]{2}
\symdef{Three}[name=three]{3}
\symdef{Four}[name=four]{4}
\symdef{Five}[name=five]{5}
\symdef{Six}[name=six]{6}
\symdef{Seven}[name=seven]{7}
\symdef{Eight}[name=eight]{8}
\symdef{Nine}[name=nine]{9}

\symdecl*{natural number}
\symdecl*{nonzero}
\symdecl*{direct successor}
\symdecl*{injectivity of successor function}
\symdecl*{induction I}
\symdecl*{right-cancellability of addition}
\symdecl*{left-cancellability of addition}

\begin{signature*}[forthel,for=natural number]
  A \emph{natural number} is an object.
\end{signature*}

\begin{definition*}[forthel,for=Naturals]
  $\emph{\Naturals}$ is the class of natural numbers.
\end{definition*}

\begin{signature*}[forthel,for=sum]
  Let $n, m$ be natural numbers.
  $\emph{n \Plus m}$ is a natural number.
  Let the \emph{sum of $n$ and $m$} stand for $n \Plus m$.
\end{signature*}

\begin{signature*}[forthel,for={zero,nonzero}]
  $\emph{\Zero}$ is a natural number.
  Let \emph{zero} stand for $\Zero$.
  Let $n$ is \emph{nonzero} stand for $n \NotEq \Zero$.
\end{signature*}

\begin{signature*}[forthel,for={one,direct successor}]
  $\emph{\One}$ is a natural number.
  Let \emph{one} stand for $\One$.
  Let the \emph{direct successor of $n$} stand for $n \Plus \One$.
\end{signature*}

\begin{definition*}[forthel,for=two]
  $\emph{\Two} \DefEq \One \Plus \One$.
  Let \emph{two} stand for $\Two$.
\end{definition*}

\begin{definition*}[forthel,for=three]
  $\emph{\Three} \DefEq \Two \Plus \One$.
  Let \emph{three} stand for $\Three$.
\end{definition*}

\begin{definition*}[forthel,for=four]
  $\emph{\Four} \DefEq \Three \Plus \One$.
  Let \emph{four} stand for $\Four$.
\end{definition*}

\begin{definition*}[forthel,for=five]
  $\emph{\Five} \DefEq \Four \Plus \One$.
  Let \emph{five} stand for $\Five$.
\end{definition*}

\begin{definition*}[forthel,for=six]
  $\emph{\Six} \DefEq \Five \Plus \One$.
  Let \emph{six} stand for $\Six$.
\end{definition*}

\begin{definition*}[forthel,for=seven]
  $\emph{\Seven} \DefEq \Six \Plus \One$.
  Let \emph{seven} stand for $\Seven$.
\end{definition*}

\begin{definition*}[forthel,for=eight]
  $\emph{\Eight} \DefEq \Seven \Plus \One$.
  Let \emph{eight} stand for $\Eight$.
\end{definition*}

\begin{definition*}[forthel,for=nine]
  $\emph{\Nine} \DefEq \Eight \Plus \One$.
  Let \emph{nine} stand for $\Nine$.
\end{definition*}

\begin{axiom*}[forthel,name=injectivity of successor function]
  Let $n, m$ be natural numbers.
  If $n \Plus \One \Eq m \Plus \One$ then $n \Eq m$.
\end{axiom*}

\begin{axiom*}[forthel]
  There exists no natural number $n$ such that $n \Plus \One \Eq \Zero$.
\end{axiom*}

\begin{axiom*}[forthel,title=Induction,name=induction I]
  Let $\Phi$ be a class.
  Assume $\Zero \Elem \Phi$ and for all natural numbers $n$ if $n \Elem \Phi$ then
  $n \Plus \One \Elem \Phi$.
  Then $\Phi$ contains every natural number.
\end{axiom*}

\begin{axiom*}[forthel]
  Then $\One \Eq \Zero \Plus \One$.
\end{axiom*}

\begin{axiom*}[forthel]
  Let $n$ be a natural number.
  Then $n \Plus \Zero \Eq n$.
\end{axiom*}

\begin{axiom*}[forthel]
  Let $n, m$ be natural numbers.
  Then $n \Plus (m \Plus \One) \Eq (n \Plus m) \Plus \One$.
\end{axiom*}

\begin{proposition*}[forthel]
  Let $n$ be a natural number.
  Then $n \Eq \Zero$ or $n \Eq m \Plus \One$ for some natural number $m$.
\end{proposition*}
\begin{proof}[forthel]
  Define $\Phi \DefEq \SClass{n'}{\Naturals}{n' \Eq \Zero\text{ or }n' \Eq m' \Plus \One\text{ for some natural number }m'}$.
  $\Zero \Elem \Phi$ and for all $n' \Elem  \Phi$ we have $n' \Plus \One \Elem \Phi$.
  Hence every natural number is contained in $\Phi$.
  Thus $n \Eq \Zero$ or $n \Eq m \Plus \One$ for some natural number $m$.
\end{proof}

\begin{proposition*}[forthel]
  Let $n$ be a natural number.
  Then $n \NotEq n \Plus \One$.
\end{proposition*}
\begin{proof}[forthel]
  Define $\Phi \DefEq \SClass{n'}{\Naturals}{n' \NotEq n' \Plus \One}$.

  (1) $\Zero$ belongs to $\Phi$.

  (2) For all $n' \Elem \Phi$ we have $n' \Plus \One \Elem \Phi$.
  \begin{proof}
    Let $n' \Elem \Phi$.
    Then $n' \NotEq n' \Plus \One$.
    If $n' \Plus \One \Eq (n' \Plus \One) \Plus \One$ then $n' \Eq n' \Plus \One$.
    Thus it is wrong that $n' \Plus \One \Eq (n' \Plus \One) \Plus \One$.
    Hence $n' \Plus \One \Elem \Phi$.
  \end{proof}

  Therefore every natural number is an element of $\Phi$.
  Consequently $n \NotEq n \Plus \One$.
\end{proof}

\begin{proposition*}[forthel]
  Let $n, m, k$ be natural numbers.
  Then $n \Plus (m \Plus k) \Eq (n \Plus m) \Plus k$.
\end{proposition*}
\begin{proof}[forthel]
  Define $\Phi \DefEq \SClass{k'}{\Naturals}{n \Plus (m \Plus k') \Eq (n \Plus m) \Plus k'}$.

  (1) $\Zero$ is contained in $\Phi$.
  Indeed $n \Plus (m \Plus \Zero) \Eq n \Plus m \Eq (n \Plus m) \Plus \Zero$.

  (2) For all $k' \Elem \Phi$ we have $k' \Plus \One \Elem \Phi$.
  \begin{proof}
    Let $k' \Elem \Phi$.
    Then $n \Plus (m \Plus k') \Eq (n \Plus m) \Plus k'$.
    Hence
    \[  n \Plus (m \Plus (k' \Plus \One))        \]
    \[    \Eq n \Plus ((m \Plus k') \Plus \One)    \]
    \[    \Eq (n \Plus (m \Plus k')) \Plus \One    \]
    \[    \Eq ((n \Plus m) \Plus k') \Plus \One    \]
    \[    \Eq (n \Plus m) \Plus (k' \Plus \One).   \]
    Thus $k' \Plus \One \Elem \Phi$.
  \end{proof}

  Thus every natural number is an element of $\Phi$.
  Therefore $n \Plus (m \Plus k) \Eq (n \Plus m) \Plus k$.
\end{proof}

\begin{proposition*}[forthel]
  Let $n, m$ be natural numbers.
  Then $n \Plus m \Eq m \Plus n$.
\end{proposition*}
\begin{proof}[forthel]
  Define $\Phi \DefEq \SClass{m'}{\Naturals}{n \Plus m' \Eq m' \Plus n}$.

  (1) $\Zero$ is an element of $\Phi$.
  \begin{proof}
    Define $\Psi \DefEq \SClass{n'}{\Naturals}{n' \Plus \Zero \Eq \Zero \Plus n'}$.

    (1a) $\Zero$ belongs to $\Psi$.

    (1b) For all $n' \Elem \Psi$ we have $n' \Plus \One \Elem \Psi$.
    \begin{proof}
      Let $n' \Elem \Psi$.
      Then $n' \Plus \Zero \Eq \Zero \Plus n'$.
      Hence
      \[  (n' \Plus \One) \Plus \Zero        \]
      \[    \Eq n' \Plus \One          \]
      \[    \Eq (n' \Plus \Zero) \Plus \One    \]
      \[    \Eq (\Zero \Plus n') \Plus \One    \]
      \[    \Eq \Zero \Plus (n' \Plus \One).   \]
    \end{proof}

    Hence every natural number belongs to $\Psi$.
    Thus $n \Plus \Zero \Eq \Zero \Plus n$.
    Therefore $\Zero$ is an element of $\Phi$.
  \end{proof}

  Let us show that (2) $n \Plus \One \Eq \One \Plus n$.
  \begin{proof}
    Define $\Theta \DefEq \SClass{n'}{\Naturals}{n' \Plus \One \Eq \One \Plus n'}$.

    (2a) $\Zero$ is an element of $\Theta$.

    (2b) For all $n' \Elem \Theta$ we have $n' \Plus \One \Elem \Theta$.
    \begin{proof}
      Let $n' \Elem \Theta$.
      Then $n' \Plus \One \Eq \One \Plus n'$.
      Hence
      \[  (n' \Plus \One) \Plus \One        \]
      \[    \Eq (\One \Plus n') \Plus \One    \]
      \[    \Eq \One \Plus (n' \Plus \One).   \]
      Thus $n' \Plus \One \Elem \Theta$.
    \end{proof}

    Thus every natural number belongs to $\Theta$.
    Therefore $n \Plus \One \Eq \One \Plus n$.
  \end{proof}

  (3) For all $m' \Elem \Phi$ we have $m' \Plus \One \Elem \Phi$.
  \begin{proof}
    Let $m' \Elem \Phi$.
    Then $n \Plus m' \Eq m' \Plus n$.
    Hence
    \[  n \Plus (m'  \Plus \One)       \]
    \[    \Eq (n \Plus m') \Plus \One    \]
    \[    \Eq (m' \Plus n) \Plus \One    \]
    \[    \Eq m' \Plus (n \Plus \One)    \]
    \[    \Eq m' \Plus (\One \Plus n)    \]
    \[    \Eq (m' \Plus \One) \Plus n.   \]
    Thus $m' \Plus \One \Elem \Phi$.
  \end{proof}

  Thus every natural number is an element of $\Phi$.
  Therefore $n \Plus m \Eq m \Plus n$.
\end{proof}

\begin{proposition*}[forthel,name=right-cancellability of addition]
  Let $n, m, k$ be natural numbers.
  If $n \Plus k \Eq m \Plus k$ then $n \Eq m$.
\end{proposition*}
\begin{proof}[forthel]
  Define $\Phi \DefEq \SClass{k'}{\Naturals}{\text{ if }n \Plus k' \Eq m \Plus k'\text{ then }n \Eq m}$.

  (1) $\Zero$ is an element of $\Phi$.

  (2) For all $k' \Elem \Phi$ we have $k' \Plus \One \Elem \Phi$.
  \begin{proof}
    Let $k' \Elem \Phi$.
    Suppose $n \Plus (k' \Plus \One) \Eq m \Plus (k' \Plus \One)$.
    Then $(n \Plus k') \Plus \One \Eq (m \Plus k') \Plus \One$.
    Hence $n \Plus k' \Eq m \Plus k'$.
    Thus $n \Eq m$.
  \end{proof}

  Therefore every natural number is an element of $\Phi$.
  Consequently if $n \Plus k \Eq m \Plus k$ then $n \Eq m$.
\end{proof}

\begin{corollary*}[forthel,name=left-cancellability of addition]
  Let $n, m, k$ be natural numbers.
  If $k \Plus n \Eq k \Plus m$ then $n \Eq m$.
\end{corollary*}
\begin{proof}[forthel]
  Assume $k \Plus n \Eq k \Plus m$.
  We have $k \Plus n \Eq n \Plus k$ and $k \Plus m \Eq m \Plus k$.
  Hence $n \Plus k \Eq m \Plus k$.
  Thus $n \Eq m$.
\end{proof}

\begin{proposition*}[forthel]
  Let $n, m$ be natural numbers.
  If $n \Plus m \Eq \Zero$ then $n \Eq \Zero$ and $m \Eq \Zero$.
\end{proposition*}
\begin{proof}[forthel]
  Assume $n \Plus m \Eq \Zero$.
  Suppose $n \NotEq \Zero$ or $m \NotEq \Zero$.
  Then we can take a $k \Elem \Naturals$ such that $n \Eq k \Plus \One$ or $m \Eq k \Plus \One$.
  Hence there exists a natural number $l$ such that
  $n \Plus m
    \Eq l \Plus (k \Plus \One)
    \Eq (l \Plus k) \Plus \One
    \NotEq \Zero$.
  Contradiction.
\end{proof}

\end{smodule}
\end{document}
